\documentclass[11pt,letterpaper]{article}
\usepackage[spanish,es-nodecimaldot]{babel}
\usepackage[utf8]{inputenc}
\usepackage[T1]{fontenc}
\usepackage{lmodern}
\usepackage{geometry}
\usepackage{array,booktabs,longtable,tabularx}
\usepackage{amssymb}
\usepackage[table]{xcolor}
\usepackage{enumitem}
\usepackage{graphicx}
\usepackage{fancyhdr}
\usepackage{hyperref}
\usepackage{microtype}
\usepackage{pdflscape}
\geometry{margin=2cm}
\hypersetup{colorlinks=true,linkcolor=blue!55!black,urlcolor=blue!55!black}
\setlength{\parindent}{0pt}
\setlength{\parskip}{5pt}
\renewcommand{\arraystretch}{1.22}
\pagestyle{fancy}
\setlength{\headheight}{14pt}
\fancyhf{}
\lhead{Taller 2 -- Corte 1}
\rhead{Nova Spot Micro}
\cfoot{\thepage}
\newcommand{\verde}{\cellcolor{green!18}\textbf{Verde}}
\newcommand{\amarillo}{\cellcolor{yellow!25}\textbf{Amarillo}}
\newcommand{\rojo}{\cellcolor{red!16}\textbf{Rojo}}
\newcommand{\estado}[1]{\textbf{#1}}

\begin{document}

\begin{titlepage}
\centering
{\Large Universidad Santo Tomás\par}
\vspace{0.5cm}
{\large Proyecto de Grado II -- Taller 2, Corte 1\par}
\vfill
{\Huge\bfseries Matriz de transición del anteproyecto al documento final\par}
\vspace{0.8cm}
{\Large Diseño de un sistema de control para la locomoción y estabilización de un robot cuadrúpedo basado en aprendizaje por refuerzo\par}
\vfill
\begin{tabular}{rl}
\textbf{Estudiantes:} & David Esteban Díaz Castro\\
& David Felipe Díaz Suesca\\[2mm]
\textbf{Director:} & MsC. Armando Mateus Rojas\\
\textbf{Codirectores:} & PhD. Billy Vladimir Toro Tovar\\
& MsC. Carlos Javier Mojica Casallas\\[2mm]
\textbf{Fecha de corte:} & 31 de agosto de 2026
\end{tabular}
\vfill
{\small Documento elaborado a partir del anteproyecto, la tesis en curso, el cronograma y la evidencia reproducible disponible en el repositorio del proyecto.}
\end{titlepage}

\tableofcontents
\newpage

\section{Propósito, alcance y estado del proyecto}

Este documento responde al Taller 2 mediante una auditoría entre lo prometido y
lo ejecutado. El proyecto conserva como objetivo final comparar una marcha
nominal con la misma marcha asistida por una capa correctiva de aprendizaje por
refuerzo (RL). A la fecha del corte no se declara cumplido el objetivo general:
la marcha convencional y el entorno de simulación están avanzados, mientras que
la caracterización definitiva, la seguridad física, RL y la comparación final
siguen pendientes.

La evidencia más madura corresponde al modelo matemático, ROS 2, Gazebo,
MuJoCo y las marchas nominales. El gateo de cadencia corregida fue reproducido
en dos ensayos de 13 ciclos, con diferencias menores al 0,2\% en las medias. La
marcha paso fue validada en dos ensayos Gazebo de 12 ciclos y en MuJoCo. El
análisis de contactos reveló, sin embargo, que caminar visualmente no implica
cumplir el patrón ideal de apoyos. Este hallazgo motivó iteraciones controladas
sin ocultar tentativas inválidas ni congelar prematuramente una nueva línea
base.

La caracterización física está temporalmente bloqueada porque falta imprimir y
montar \path{SM3_Cover_LeftFemur.stl}. Por ello no se presentan dimensiones,
masas o calibraciones parciales como propiedades definitivas del robot.

\section{A. Auditoría de secciones heredadas}

\subsection{Hallazgos iniciales}

\begin{tabularx}{\textwidth}{>{\bfseries}p{3.5cm}X}
\toprule
Aspecto que ya no representa con precisión el proyecto & El texto heredado sugiere por momentos un modelo dinámico completo o una reactivación física cercana. La evidencia actual sustenta un modelo nominal computable y simulación reproducible, no un gemelo digital identificado ni una plataforma física totalmente validada.\\
Aspecto para revisión con el director & Métrica primaria de la comparación, umbral de mejora RL, límite de continuidad articular, criterio final de caída, alcance del modelo dinámico y permanencia del requisito HRI.\\
\bottomrule
\end{tabularx}

\begin{landscape}
\small
\begin{longtable}{p{2.6cm}p{5.0cm}p{6.1cm}p{1.8cm}p{6.0cm}}
\caption{Auditoría de las secciones del anteproyecto.}\\
\toprule
\textbf{Sección} & \textbf{Qué establece} & \textbf{Qué ocurrió durante el desarrollo} & \textbf{Estado} & \textbf{Acción requerida}\\
\midrule
\endfirsthead
\toprule
\textbf{Sección} & \textbf{Qué establece} & \textbf{Qué ocurrió durante el desarrollo} & \textbf{Estado} & \textbf{Acción requerida}\\
\midrule
\endhead
Problema & Reactivar una plataforma inactiva mediante caracterización, electrónica, modelado y control. & Se construyeron modelos y control nominal reproducible en simulación; la plataforma física continúa incompleta y no se ha ejecutado la comparación RL. & A & Distinguir resultados demostrados, trabajo pendiente y proyección social futura.\\
Justificación & Recuperación académica de la plataforma y posibles aplicaciones en entornos peligrosos. & La contribución formativa y técnica se sostiene. La inspección peligrosa sigue siendo una aplicación futura, no validada. & F & Fortalecer con resultados reales y mantener explícita la diferencia entre aporte logrado y proyección.\\
Objetivo general & Combinar marcha nominal y RL correctivo para mejorar estabilidad y repetibilidad. & La marcha nominal existe; RL y comparación final aún no se han ejecutado. & C & Conservar sin afirmar cumplimiento anticipado.\\
OE1: caracterización y modelos & Caracterizar y desarrollar modelos cinemático y dinámico. & FK, IK, Jacobiano, dinámica nominal, URDF y MJCF implementados. Faltan geometría, masas y parámetros físicos definitivos. & A & Declarar modelo nominal y completar tabla nominal--medido tras el ensamble.\\
OE2: arquitectura segura & Diseñar electrónica y software con límites eléctricos, mecánicos y articulares. & ROS 2, supervisor provisional, Raspberry, Mega, I2C y PCA9685 verificados; faltan OE con pull-up, parada física, fuente bajo carga y 12 calibraciones. & F & Incorporar evidencias eléctricas, curvas de calibración y pruebas provocadas del supervisor.\\
OE3: locomoción convencional & Implementar FSM para postura, paso y gateo. & Postura, gateo y marcha paso se ejecutan en Gazebo; paso también fue verificado articularmente en MuJoCo. & F & Corregir contacto trasero, congelar línea base y validar posteriormente en hardware.\\
OE4: política RL & Entrenar e integrar correcciones pequeñas, sin PWM directo. & Semillas y alcance definidos, pero PPO permanece bloqueado hasta cerrar marcha física y protocolo. & C & Mantener el objetivo y documentar observaciones, acciones, recompensa y saturaciones antes de entrenar.\\
OE5: comparación & Comparar nominal y nominal+RL en simulación y robot físico. & Existe matriz de cinco ensayos por condición y marcha; todavía no hay políticas ni 400 ciclos finales. & C & Ejecutar solamente después de congelar protocolo, política y línea base.\\
Alcance & Superficie plana, baja velocidad, paso y gateo; incluye HRI y modelo dinámico, excluye navegación y galope rápido. & Se mantuvo superficie plana y control acotado. El galope fue relegado formalmente. HRI no tiene aún producto verificable. Hay contradicción sobre la complejidad dinámica. & R & Acordar con el director la redacción del modelo nominal y decidir si HRI permanece como entregable.\\
Marco referencial & Cinemática, dinámica, marchas, control y RL. & Los conceptos se usan en código y simulación; deben conectarse mejor con decisiones, métricas y limitaciones. & F & Eliminar teoría desconectada e incorporar referencias para PPO, sim-to-real, contacto y seguridad.\\
Metodología & Ocho fases secuenciales con retroalimentación y transferencia progresiva. & La simulación se adelantó a la caracterización física debido al estado incompleto del robot. Se incorporaron rosbag2, hashes, fases y contactos como evidencia. & A & Reescribir como metodología realmente ejecutada y justificar los ciclos de corrección.\\
\bottomrule
\end{longtable}
\end{landscape}

\begin{landscape}
\section{B. Matriz objetivo--evidencia}
\scriptsize
\begin{longtable}{p{1.7cm}p{2.8cm}p{3.5cm}p{4.2cm}p{3.5cm}p{3.5cm}p{1.8cm}}
\caption{Correspondencia entre objetivos y evidencia al 31 de agosto de 2026.}\\
\toprule
\textbf{Objetivo} & \textbf{Qué prometía} & \textbf{Qué se hizo} & \textbf{Evidencia disponible} & \textbf{Evidencia faltante} & \textbf{Métrica o criterio} & \textbf{Clasificación}\\
\midrule
\endfirsthead
\toprule
\textbf{Objetivo} & \textbf{Qué prometía} & \textbf{Qué se hizo} & \textbf{Evidencia disponible} & \textbf{Evidencia faltante} & \textbf{Métrica o criterio} & \textbf{Clasificación}\\
\midrule
\endhead
OE1 & Caracterización y modelos para simulación y referencias. & Modelo matemático propio, FK/IK, Jacobiano, dinámica nominal, URDF/MJCF y pruebas. & \path{Documentacion/MODELO_MATEMATICO_LATEX}; pruebas de cinemática y coherencia; modelos en \path{src}. & Tres mediciones físicas, masas, holguras, identificación de servos y actualización nominal--medido. & Error FK--IK; alcanzabilidad; límites; coherencia entre modelos; incertidumbre física. & \amarillo\\
OE2 & Arquitectura electrónica y de software segura. & ROS 2, supervisor, Raspberry Ubuntu 22.04, SSH/DDS, Mega--PCA9685 y movimiento limitado CH5--CH10. & Diagramas, prueba I2C 0x40, compilación ARM64, 41 pruebas del controlador y documentación de seguridad. & OE con pull-up, parada física, fuente bajo carga, corriente/temperatura, mapa de canales y calibración de 12 servos. & Arranque deshabilitado; límites PWM; pérdida de comunicación a estado seguro; corriente y temperatura. & \amarillo\\
OE3 & Postura, paso y gateo nominales mediante FSM. & Gateo cartesiano, marcha paso, temporizador corregido, fase y contactos previstos/medidos. & Dos líneas base de gateo; dos ensayos paso; validación MuJoCo; bolsas, CSV, gráficas y hashes. & Contacto trasero coherente, línea base final congelada y validación física progresiva. & Ciclos; avance; velocidad; roll/pitch; lateral; cadencia; continuidad; fallos; contactos. & \amarillo\\
OE4 & PPO correctivo acotado e integración progresiva. & Se fijaron semillas y principio de corrección sobre referencias; no se inició entrenamiento. & Decisión de semillas 11, 23, 37, 53 y 71; supervisor independiente; galope excluido. & MDP, recompensa, saturaciones, currículo, modelos entrenados, validación separada y transferencia. & Recompensa; estabilidad; saturaciones; tasa de fallos; comparación con nominal. & \rojo\\
OE5 & Comparación experimental nominal frente a nominal+RL. & Se definieron cinco ensayos de 20 ciclos por marcha y condición, con escenarios emparejados. & Protocolo borrador y matriz de 400 ciclos; analizadores reproducibles. & Métrica primaria aprobada, políticas, ensayos finales en simulación y hardware, análisis estadístico. & Unidad: ensayo (n=5); IC; diferencias emparejadas; fallos e intervenciones. & \rojo\\
\bottomrule
\end{longtable}
\end{landscape}

\textbf{Dictamen global: parcialmente demostrable.} Un jurado puede reproducir
y evaluar OE1 y la parte simulada de OE2--OE3. No puede afirmar aún el
cumplimiento del objetivo general, OE4 u OE5.

\begin{landscape}
\section{C. Metodología proyectada frente a metodología ejecutada}
\small
\begin{longtable}{p{2.4cm}p{4.2cm}p{5cm}p{5cm}p{4cm}}
\toprule
\textbf{Aspecto} & \textbf{Lo planteado} & \textbf{Lo ejecutado} & \textbf{Cambio y justificación} & \textbf{Evidencia para la tesis}\\
\midrule
Requisitos & Reactivación, paso, gateo, RL, seguridad y HRI. & Se precisaron baja velocidad, superficie plana, RL correctivo y galope fuera del alcance. & Se redujo ambigüedad y se protegió el hardware. & Tabla de requisitos y trazabilidad.\\
Arquitectura & Control nominal más RL. & ROS 2 con controlador, métricas, supervisor, fase, contactos y dos simuladores. & Se añadieron observabilidad y procesos separados para evitar pérdida de fases. & Diagrama de nodos, tópicos y flujo de seguridad.\\
Tecnologías & Plataforma física y simulación. & ROS 2 Humble, Gazebo Sim, MuJoCo, Raspberry Pi 4, Mega y PCA9685. & Ubuntu 24.04 fue descartado para mantener compatibilidad con Humble 22.04. & Matriz de selección y versiones.\\
Modelado & Cinemática y modelo dinámico. & Modelo nominal computable con contacto simplificado. & No se denomina gemelo digital porque faltan parámetros identificados. & Ecuaciones, supuestos y tabla nominal--medido futura.\\
Cálculos & Trayectorias y estabilidad. & FK/IK, Jacobiano, masa, Coriolis, gravedad, contacto y margen estático. & Se implementaron pruebas automatizadas y ejemplos reproducibles. & Documento matemático y pruebas.\\
Simulación & Validación previa al hardware. & Ensayos con rosbag2, repetición, hashes, contactos y análisis por ciclo. & Se corrigió un temporizador que producía ciclos de 4,80 s en vez de 4,32 s. & Bolsas separadas antes/después y validación temporal.\\
Implementación & FSM para marchas. & Gateo de 24 referencias y marcha paso de 32; galope deshabilitado. & Los contactos mostraron deslizamiento trasero aunque el robot avanzaba. & Código, configuración, pruebas y diagnósticos.\\
Integración & Progresión segura al robot. & Red, DDS, I2C y servos limitados; no hay marchas físicas. & El ensamble incompleto y la seguridad pendiente impiden avanzar responsablemente. & Checklists, pieza pendiente y pruebas eléctricas.\\
Procedimiento experimental & Al menos diez ciclos por condición. & Se fijaron cinco ensayos de 20 ciclos; ciclos 2--20 en régimen y ensayo como unidad independiente. & Se evitó pseudorreplicación y selección favorable. & Protocolo y decisión de semillas.\\
Validación & Simulación y robot físico con/sin RL. & Líneas base nominales reproducibles; comparación final pendiente. & La validación se escalonó por riesgo y madurez. & Matriz objetivo--evidencia y plan inmediato.\\
\bottomrule
\end{longtable}
\end{landscape}

La secuencia documentada puede reconstruirse como
\textbf{decisión $\rightarrow$ criterio $\rightarrow$ evidencia $\rightarrow$
consecuencia}. Por ejemplo: se detectó una cadencia observada de 4,80 s;
se corrigió el avance de la fecha límite del temporizador; seis ciclos midieron
4,320106573 s; luego se generó y reprodujo una línea base nueva sin mezclarla
con bolsas históricas.

\section{D. Evidencia crítica de validación}

Se selecciona OE3 porque una marcha nominal reproducible es condición previa
para hardware, RL y comparación final.

\begin{tabularx}{\textwidth}{>{\bfseries}p{5.2cm}X}
\toprule
Elemento & Definición para el proyecto\\
\midrule
Objetivo asociado & OE3: implementar postura, marcha paso y gateo nominales mediante FSM.\\
Variable a evaluar & Cadencia, ciclos completados, pose corporal, continuidad articular, contacto previsto/observado y activaciones del supervisor.\\
Referencia esperada & Gateo: 24 referencias, 0,18 s por referencia y 4,32 s por ciclo; paso de 18 mm y elevación de 14 mm.\\
Herramienta & Gazebo Sim, ROS 2 Humble, rosbag2, fase \path{/nova/gait_phase}, contactos y analizadores Python.\\
Condiciones & Superficie plana, estado inicial equivalente, un nodo por proceso, parámetros congelados, órdenes gateo y stand registradas.\\
Repeticiones & Desarrollo: dos ensayos independientes de al menos 10 ciclos. Comparación final: cinco ensayos de 20 ciclos por condición.\\
Métricas & Duración por ciclo, avance, velocidad, excursión lateral, altura, roll/pitch, salto articular, coincidencia de contactos, fallos e intervenciones.\\
Origen del criterio & Configuración implementada, límites del supervisor y protocolo; umbral articular definitivo requiere aval del director.\\
Criterio & Todos los ciclos previstos, cero caída/intervención, cadencia dentro de $\pm1\%$, datos esenciales completos y continuidad bajo límite aprobado.\\
Resultado a reportar & Media y dispersión por ensayo, transitorio separado, fallos conservados, hashes y diferencias entre repeticiones.\\
Figura/tabla & Series temporales, resumen por ciclo, transición de contactos y tabla comparativa entre ensayos.\\
\bottomrule
\end{tabularx}

\subsection{Resultado disponible de la evidencia crítica}

La línea base de cadencia corregida y su repetición completaron 13 ciclos cada
una, sin activaciones del supervisor. Las diferencias entre medias fueron
menores al 0,2\%. En ciclos 2--13, la primera ejecución obtuvo 0,023338 m/ciclo,
0,005402 m/s, 0,015133 m de excursión lateral, roll máximo de 2,233112 grados y
pitch máximo de 4,365442 grados. Esto demuestra reproducibilidad en simulación
de la referencia articular, pero no valida todavía la geometría física ni el
patrón ideal de contacto.

\section{E. Proyección del documento final}

\subsection{Capítulo de desarrollo}

\begin{longtable}{p{6.2cm}p{9.2cm}}
\toprule
\textbf{Pregunta de ingeniería} & \textbf{Evidencia o sección requerida}\\
\midrule
¿Qué debía cumplir? & Requisitos funcionales, seguridad, superficie plana, modos y corrección RL acotada.\\
¿Qué restricciones existían? & MG996R, geometría provisional, potencia, sensado, ensamble incompleto y límites de simulación.\\
¿Qué arquitectura se definió? & Diagramas ROS 2, electrónica, potencia, comunicaciones y supervisor.\\
¿Qué alternativas se consideraron? & Gazebo/MuJoCo, perfiles de descenso, Raspberry/Mega y exclusión del galope.\\
¿Cómo se seleccionó? & Criterios de compatibilidad, continuidad, seguridad, reproducibilidad y alcance.\\
¿Qué cálculos sustentan el diseño? & FK/IK, Jacobianos, dinámica nominal, contacto, estabilidad y conversión articulación--PWM futura.\\
¿Qué simulaciones se realizaron? & Línea base, repetición, marcha paso, contactos, temporizador y curvas de liberación.\\
¿Cómo se implementó e integró? & Paquetes, nodos, parámetros, tópicos, Raspberry, DDS, I2C y PCA9685.\\
¿Qué dificultades aparecieron? & Nodos duplicados, pérdida de marcadores/fases, cadencia tardía, deslizamiento trasero y carpeta vacía en Git.\\
¿Qué modificaciones fueron necesarias? & Temporizador, QoS, separación de procesos, subfases, curvas por eje y bloqueos de seguridad.\\
¿Cómo se verificó? & Pytest, colcon, bolsas, hashes, CSV, gráficas, pruebas provocadas y repetición independiente.\\
¿Cómo se validará el sistema completo? & Comparación emparejada nominal/RL en simulación y hardware bajo protocolo aprobado.\\
\bottomrule
\end{longtable}

\subsection{Proyección de resultados y discusión}

\textbf{Objetivo seleccionado:} OE3, marcha nominal.

\textbf{Condiciones:} superficie plana, Gazebo, 24 referencias, paso de 18 mm,
elevación de 14 mm, 0,18 s por muestra y estados iniciales controlados.

\textbf{Evidencias:} tablas, series temporales, simulación, comparación entre
ensayos, CSV por ciclo y transiciones de contacto.

\textbf{Métrica principal provisional:} proporción de ensayos que completan
todos los ciclos sin fallo ni intervención. Métricas descriptivas: avance,
velocidad, roll/pitch, lateral, continuidad y contactos.

\textbf{Discusión prevista:} la cadencia y el avance fueron altamente
repetibles, pero el contacto físico simulado no coincidió con el plan. El
hallazgo muestra que estabilidad visual, avance y sincronía de apoyos son
dimensiones distintas. La discusión deberá relacionar este comportamiento con
la dinámica nominal, el timeout del sensor, la transferencia de peso y la forma
de la trayectoria, sin atribuirlo automáticamente al hardware.

\textbf{Juicio actual:} \estado{cumple parcialmente}. Existe una marcha nominal
reproducible en simulación, pero faltan contacto corregido, congelación final y
validación física.

\begin{landscape}
\section{Conclusiones preliminares sustentables}
\small
\begin{longtable}{p{1.6cm}p{4.5cm}p{4.5cm}p{4.7cm}p{3.8cm}}
\toprule
\textbf{Objetivo} & \textbf{Resultado principal} & \textbf{Evidencia cuantitativa} & \textbf{Conclusión sustentable hoy} & \textbf{Limitación}\\
\midrule
OE1 & Modelo nominal y simuladores construidos. & Pruebas FK/IK, 41 pruebas del paquete y documentos matemáticos. & Es posible generar referencias alcanzables y coherentes para el modelo provisional. & Parámetros físicos no identificados.\\
OE2 & Arquitectura lógica y comunicación verificadas. & DDS recibió ocho mensajes consecutivos; PCA9685 detectado en 0x40; CH5--CH10 movidos en rango limitado. & La cadena de comunicación funciona por subsistemas. & Seguridad eléctrica y 12 calibraciones incompletas.\\
OE3 & Gateo y paso repetibles en simulación. & Gateo: diferencias menores al 0,2\%; paso: menores al 0,4\%; cero eventos en ensayos reportados. & La FSM produce locomoción nominal reproducible en simulación bajo condiciones congeladas. & Contactos y hardware pendientes.\\
OE4 & No hay resultado de política. & Solo semillas y restricciones congeladas. & No puede concluirse mejora por RL. & Falta todo el entrenamiento y validación.\\
OE5 & Diseño comparativo preparado. & 5 ensayos $\times$ 20 ciclos $\times$ 2 condiciones $\times$ 2 marchas = 400 ciclos. & Existe un protocolo planificado, no evidencia comparativa. & Métrica primaria y ejecución pendientes.\\
\bottomrule
\end{longtable}
\end{landscape}

Conclusión de una frase para OE3: \emph{En Gazebo y con la configuración
congelada, el controlador convencional ejecutó gateo de 4,32 s por ciclo en dos
ensayos independientes de 13 ciclos, sin eventos del supervisor y con
diferencias menores al 0,2\% entre medias, lo que demuestra reproducibilidad de
la locomoción simulada, aunque no valida todavía los contactos ni el robot
físico.}

\section{Revisión tipo jurado técnico}

La autoevaluación asigna el dictamen \textbf{PARCIALMENTE DEMOSTRADO}. La
evidencia tiene ubicación, variables, métricas, condiciones y procedimientos
reproducibles para la marcha simulada. El criterio articular definitivo no está
aprobado y faltan caracterización física, seguridad, RL y comparación final.
La revisión entre pares real debe añadirse después de recibir el nombre y las
observaciones del grupo revisor; no se inventa en este informe.

\section{Plan inmediato de trabajo}

Prioridades: (1) protocolo y criterios con el director; (2) contacto y línea
base nominal; (3) completar ensamble, caracterización y seguridad antes de
marchas físicas.

\small
\begin{longtable}{p{5.1cm}p{4.4cm}p{1.5cm}p{2.4cm}p{2.2cm}}
\toprule
\textbf{Acción} & \textbf{Evidencia esperada} & \textbf{Pri.} & \textbf{Responsable} & \textbf{Fecha prevista}\\
\midrule
Aprobar métrica primaria, frecuencia y umbrales. & Ficha firmada y protocolo v2. & Alta & Equipo/director & 07/09/2026\\
Revisar semántica y debounce de contactos traseros. & Prueba unitaria, criba y ensayo equivalente. & Alta & Estudiantes & 11/09/2026\\
Imprimir y montar tapa izquierda de fémur. & Pieza instalada y fotografía. & Alta & Estudiantes & Según impresión\\
Caracterizar geometría, masas y holguras. & Ficha completa, fotos, hashes y tabla nominal--medido. & Alta & Estudiantes & 7 días tras montaje\\
Cerrar OE, pull-up, parada y fuente bajo carga. & Checklist y mediciones eléctricas. & Alta & Equipo/director & 18/09/2026\\
Calibrar 12 MG996R individualmente. & YAML/CSV con centro, límites, sentido, corriente y temperatura. & Alta & Estudiantes & 25/09/2026\\
Fortalecer capítulos de alcance y metodología. & Texto coherente con modelo nominal y ejecución real. & Media & Equipo/director & 18/09/2026\\
Congelar especificación RL antes de entrenar. & MDP, recompensa, acciones, saturaciones y semillas. & Media & Estudiantes & Después de marcha física\\
\bottomrule
\end{longtable}
\normalsize

\section{Lista de verificación de calidad}

\begin{tabularx}{\textwidth}{Xccc}
\toprule
\textbf{Criterio} & \textbf{Sí} & \textbf{Parcial} & \textbf{No}\\
\midrule
Problema coherente con el proyecto & $\square$ & $\boxtimes$ & $\square$\\
Justificación sin promesas no demostradas & $\boxtimes$ & $\square$ & $\square$\\
Cada objetivo posee evidencia asociada & $\square$ & $\boxtimes$ & $\square$\\
Las evidencias permiten evaluar & $\square$ & $\boxtimes$ & $\square$\\
Métricas para resultados principales & $\boxtimes$ & $\square$ & $\square$\\
Criterios de aceptación definidos & $\square$ & $\boxtimes$ & $\square$\\
Metodología describe lo ejecutado & $\square$ & $\boxtimes$ & $\square$\\
Cambios metodológicos justificados & $\boxtimes$ & $\square$ & $\square$\\
Pruebas reproducibles & $\square$ & $\boxtimes$ & $\square$\\
Resultados relacionados con objetivos & $\boxtimes$ & $\square$ & $\square$\\
Proyección clara de resultados/discusión & $\boxtimes$ & $\square$ & $\square$\\
Conclusiones futuras sustentables & $\square$ & $\boxtimes$ & $\square$\\
Limitaciones técnicas identificadas & $\boxtimes$ & $\square$ & $\square$\\
Asuntos para revisión con director & $\boxtimes$ & $\square$ & $\square$\\
\bottomrule
\end{tabularx}

\section{Ubicación de la evidencia principal}

\begin{itemize}[leftmargin=*]
\item Modelo: \path{Documentacion/MODELO_MATEMATICO_LATEX}.
\item Protocolo: \path{Documentacion/PROTOCOLO_EXPERIMENTAL_BORRADOR.md}.
\item Gateo corregido: \path{Experimentos/rosbag2/linea_base_cadencia_corregida_20260814_1049}.
\item Repetición: \path{Experimentos/rosbag2/repeticion_cadencia_corregida_20260814_1100}.
\item Marcha paso: \path{Documentacion/MARCHA_PASO_VALIDACION.md}.
\item Contactos: \path{Experimentos/analisis/contactos_gateo_validado_20260814_1410}.
\item Última exploración: subdirectorio \path{liberacion_trasera_f020_r075_l080_valida_20260831} dentro de \path{Experimentos/analisis}.
\item Cronograma y ruta: \path{Github/CRONOGRAMA_SEMANAL.md} y \path{Github/RUTA_TRABAJO.md}.
\item Memoria de continuidad: \path{CONTINUIDAD.md}.
\end{itemize}

\section*{Cierre}

El proyecto cuenta con una base técnica reproducible y una marcha convencional
madura en simulación, pero la promesa completa del anteproyecto sigue abierta.
La prioridad académica es conservar esta distinción: mostrar con cifras lo ya
demostrado, declarar las limitaciones y ejecutar las fases físicas y RL solo
cuando sus dependencias estén cerradas.

\end{document}
